% GENERATED FILE. Edit canonical lessons, then run npm run generate:books.
% canonical-source-hash: fnv1a64:05cf683b5a517254
% canonical-lessons: MW-C258-chinta, MW-C258-bharoso, MW-C258-dhiraj, MW-C258-sahno, MW-C258-hiyo

\chapter{Worry and Patience}
\label{ch:mw-worry-and-patience}
\hlchaptermodality{258}

\begin{chapteropening}
\textbf{By the end of this chapter:} \emph{I can say what worries me and whom I trust.}

Complete the last lesson of chapter 258: I can say what worries me and whom I trust.
\end{chapteropening}

\section[\texorpdfstring{\mw{चिंता} (chintā)}{चिंता (chintā)}]{\mw{चिंता} (chintā) — worry}

\label{lesson:MW-C258-chinta}



\begin{quote}
Before the new one: say the Marwadi for a sigh, then the Marwadi for a reproach.
\end{quote}

\subsection*{You'll want to know: \mw{चिंता}}
\textbf{\mw{चिंता}} — \emph{chintā} — \textquotedblleft{}worry\textquotedblright{}.

\begin{quote}
\mw{चिंता} \mw{मत} \mw{करो।}
\end{quote}

\begin{quote}
\emph{chintā mat karo.}
\end{quote}

\begin{quote}
Don't worry.
\end{quote}

A feminine noun.

\begin{cousinweb}
\textbf{\mw{चिंता}} (\emph{chintā}) is a tatsama, the Sanskrit \emph{cintā}, \textquotedblleft{}thought, care\textquotedblright{}, taken over unchanged (CDIAL 4816).
\end{cousinweb}

\begin{grammarlens}[title={What you've built}]
One more word for this chapter.
\end{grammarlens}

\subsection*{Your turn}
\begin{itemize}
\raggedright
  \item \emph{Say it:} \emph{chintā}
  \item \emph{Say it:} \emph{chintā}, once more
  \item \emph{Say it:} \emph{chintā mat karo.} — Don't worry.
  \item \emph{Recall:} say the Marwadi for a sigh, then the Marwadi for a reproach, then say \emph{chintā} again
\end{itemize}

\begin{tcolorbox}[breakable,colback=teal!4,colframe=teal!35!black,title={Before you move on}]
What does \mw{चिंता} mean? (\textquotedblleft{}Worry\textquotedblright{}.) Say it once more.
\end{tcolorbox}
\section[\texorpdfstring{\mw{भरोसो} (bharoso)}{भरोसो (bharoso)}]{\mw{भरोसो} (bharoso) — trust; confidence}

\label{lesson:MW-C258-bharoso}



\begin{quote}
Before the new one: say the Marwadi for a reproach, then the Marwadi for worry.
\end{quote}

\subsection*{You'll want to know: \mw{भरोसो}}
\textbf{\mw{भरोसो}} — \emph{bharoso} — \textquotedblleft{}trust; confidence\textquotedblright{}.

\begin{quote}
\mw{भाई} \mw{माथै} \mw{म्हारो} \mw{पूरो} \mw{भरोसो} \mw{है।}
\end{quote}

\begin{quote}
\emph{bhāī māthai mhāro pūro bharoso hai.}
\end{quote}

\begin{quote}
I trust my brother completely.
\end{quote}

A masculine noun.

\begin{cousinweb}
\textbf{\mw{भरोसो}} (\emph{bharoso}) is inherited from a reconstructed \emph{*bharavaśya}, \textquotedblleft{}dependence on support\textquotedblright{}; CDIAL lists Marwari \emph{bharoso}, \textquotedblleft{}confidence\textquotedblright{}, under it (CDIAL 9398).
\end{cousinweb}

\begin{grammarlens}[title={What you've built}]
One more word for this chapter.
\end{grammarlens}

\subsection*{Your turn}
\begin{itemize}
\raggedright
  \item \emph{Say it:} \emph{bharoso}
  \item \emph{Say it:} \emph{bharoso}, once more
  \item \emph{Say it:} \emph{bhāī māthai mhāro pūro bharoso hai.} — I trust my brother completely.
  \item \emph{Recall:} say the Marwadi for a reproach, then the Marwadi for worry, then say \emph{bharoso} again
\end{itemize}

\begin{tcolorbox}[breakable,colback=teal!4,colframe=teal!35!black,title={Before you move on}]
What does \mw{भरोसो} mean? (\textquotedblleft{}Trust; confidence\textquotedblright{}.) Say it once more.
\end{tcolorbox}
\section[\texorpdfstring{\mw{धीरज} (dhīraj)}{धीरज (dhīraj)}]{\mw{धीरज} (dhīraj) — patience}

\label{lesson:MW-C258-dhiraj}



\begin{quote}
Before the new one: say the Marwadi for worry, then the Marwadi for trust.
\end{quote}

\subsection*{You'll want to know: \mw{धीरज}}
\textbf{\mw{धीरज}} — \emph{dhīraj} — \textquotedblleft{}patience\textquotedblright{}.

\begin{quote}
\mw{धीरज} \mw{राखो।}
\end{quote}

\begin{quote}
\emph{dhīraj rākho.}
\end{quote}

\begin{quote}
Be patient.
\end{quote}

A masculine noun.

\begin{cousinweb}
\textbf{\mw{धीरज}} (\emph{dhīraj}) goes back to Sanskrit \emph{dhairya}, \textquotedblleft{}firmness, patience\textquotedblright{} (Monier-Williams), in a spoken shape.
\end{cousinweb}

\begin{grammarlens}[title={What you've built}]
One more word for this chapter.
\end{grammarlens}

\subsection*{Your turn}
\begin{itemize}
\raggedright
  \item \emph{Say it:} \emph{dhīraj}
  \item \emph{Say it:} \emph{dhīraj}, once more
  \item \emph{Say it:} \emph{dhīraj rākho.} — Be patient.
  \item \emph{Recall:} say the Marwadi for worry, then the Marwadi for trust, then say \emph{dhīraj} again
\end{itemize}

\begin{tcolorbox}[breakable,colback=teal!4,colframe=teal!35!black,title={Before you move on}]
What does \mw{धीरज} mean? (\textquotedblleft{}Patience\textquotedblright{}.) Say it once more.
\end{tcolorbox}
\section[\texorpdfstring{\mw{सहणो} (sahṇo)}{सहणो (sahṇo)}]{\mw{सहणो} (sahṇo) — to bear; to put up with}

\label{lesson:MW-C258-sahno}



\begin{quote}
Before the new one: say the Marwadi for trust, then the Marwadi for patience.
\end{quote}

\subsection*{You'll want to know: \mw{सहणो}}
\textbf{\mw{सहणो}} — \emph{sahṇo} — \textquotedblleft{}to bear; to put up with\textquotedblright{}.

\begin{quote}
\mw{हूं} \mw{दरद} \mw{सहूं} \mw{हूं।}
\end{quote}

\begin{quote}
\emph{hū̃ darad sahū̃ hū̃.}
\end{quote}

\begin{quote}
I bear the pain.
\end{quote}

Past (I …): \mw{म्हैं} \mw{दरद} \mw{सह्यो}, \emph{mhaĩ darad sahyo}.

\begin{cousinweb}
\textbf{\mw{सहणो}} (\emph{sahṇo}) is inherited from Sanskrit \emph{sahate}, \textquotedblleft{}bears, suffers\textquotedblright{}; CDIAL lists Old Marwari \emph{sahaï}, \textquotedblleft{}bears\textquotedblright{}, under it (CDIAL 13304).
\end{cousinweb}

\begin{grammarlens}[title={What you've built}]
One more word for this chapter.
\end{grammarlens}

\subsection*{Your turn}
\begin{itemize}
\raggedright
  \item \emph{Say it:} \emph{sahṇo}
  \item \emph{Say it:} \emph{sahṇo}, once more
  \item \emph{Say it:} \emph{hū̃ darad sahū̃ hū̃.} — I bear the pain.
  \item \emph{Recall:} say the Marwadi for trust, then the Marwadi for patience, then say \emph{sahṇo} again
\end{itemize}

\begin{tcolorbox}[breakable,colback=teal!4,colframe=teal!35!black,title={Before you move on}]
What does \mw{सहणो} mean? (\textquotedblleft{}To bear; to put up with\textquotedblright{}.) Say it once more.
\end{tcolorbox}
\section[\texorpdfstring{\mw{हियो} (hiyo)}{हियो (hiyo)}]{\mw{हियो} (hiyo) — the heart (as the seat of feeling)}

\label{lesson:MW-C258-hiyo}



\begin{quote}
Before the new one: say the Marwadi for patience, then the Marwadi for to bear (pain).
\end{quote}

\subsection*{You'll want to know: \mw{हियो}}
\textbf{\mw{हियो}} — \emph{hiyo} — \textquotedblleft{}the heart (as the seat of feeling)\textquotedblright{}.

\begin{quote}
\mw{आज} \mw{म्हारो} \mw{हियो} \mw{राजी} \mw{है।}
\end{quote}

\begin{quote}
\emph{āj mhāro hiyo rājī hai.}
\end{quote}

\begin{quote}
My heart is glad today.
\end{quote}

A masculine noun.

\begin{cousinweb}
\textbf{\mw{हियो}} (\emph{hiyo}) is inherited from Sanskrit \emph{hṛdaya}, \textquotedblleft{}heart\textquotedblright{}; CDIAL lists Marwari \emph{hīyo} under it (CDIAL 14152).
\end{cousinweb}

\begin{grammarlens}[title={What you've built}]
One more word for this chapter.
\end{grammarlens}

\subsection*{Your turn}
\begin{itemize}
\raggedright
  \item \emph{Say it:} \emph{hiyo}
  \item \emph{Say it:} \emph{hiyo}, once more
  \item \emph{Say it:} \emph{āj mhāro hiyo rājī hai.} — My heart is glad today.
  \item \emph{Recall:} say the Marwadi for patience, then the Marwadi for to bear (pain), then say \emph{hiyo} again
\end{itemize}

\begin{tcolorbox}[breakable,colback=teal!4,colframe=teal!35!black,title={Before you move on}]
What does \mw{हियो} mean? (\textquotedblleft{}The heart (as the seat of feeling)\textquotedblright{}.) Say it once more.
\end{tcolorbox}
